% 第 2 章：SolverData 装配与网络矩阵
% 本章文件由 \input 引入，不含导言区。
\section{SolverData 装配与网络矩阵}\label{sec:assembly}

\subsection{SolverData 结构}

\compmeta{SolverData}{\srcpath{include/hacdcpf/assembly/solver\_data.hpp}}

\subsubsection*{功能说明}
\texttt{SolverData} 是潮流求解的\textbf{扁平 POD 数据底座}：rich 工程模型经
canonical 投影后，全部网络、设备、预聚合注入与矩阵以 0-based 向量形式集中存放，
求解器不再接触工程语义层。它同时携带投影回溯映射与精确 Newton 的开关量，
是"装配一次、多求解器复用"的核心（FDPF、DC、Newton、CPF 校正器均以其为输入）。

\subsubsection*{字段分组}
（\srcpath{include/hacdcpf/assembly/solver\_data.hpp:SolverData}）
\begin{itemize}
  \item \textbf{网络}：\fld{ac\_buses}、\fld{ac\_branches}、\fld{dc\_buses}、
        \fld{dc\_branches}、\fld{converters}（VSC）、\fld{dcdc\_converters}、
        \fld{energy\_routers}、\fld{switches}；
  \item \textbf{设备}：AC 侧 \fld{generators}、\fld{loads}、\fld{flexible\_loads}、
        \fld{static\_generators}、\fld{renewable\_gens}、\fld{pv\_systems}、
        \fld{storage\_units}；DC 侧 \fld{dc\_loads}、\fld{dc\_storage}、
        \fld{dc\_static\_generators}、\fld{dc\_pv\_arrays}；另有 \fld{shunts}、
        \fld{charging\_stations}、\fld{chargers}、\fld{external\_grids}、
        \fld{vpps}、\fld{microgrids}、\fld{mobile\_storage}；
  \item \textbf{标量与 ZIP}：\fld{base\_mva}=100、\fld{loss\_model}=Linear、
        \fld{zip\_pw[3]}/\fld{zip\_qw[3]}=\{1,0,0\}；
  \item \textbf{负荷聚合缓存}：\fld{has\_component\_loads}、\fld{pd\_pu}/\fld{qd\_pu}、
        每母线 ZIP 权重 \fld{bus\_zip\_zp}/\fld{\_ip}/\fld{\_pp}/\fld{\_zq}/\fld{\_iq}/\fld{\_pq}；
  \item \textbf{精确 Newton 开关}：\fld{enable\_coupled\_jacobian}、
        \fld{enable\_augmented\_equations}、\fld{enable\_semi\_smooth\_newton}、\fld{ncp\_mu}；
  \item \textbf{预聚合注入与矩阵}：\fld{pg}、\fld{qg}（Eigen 向量）、
        \fld{ybus}（稀疏复数阵）、\fld{gdc}（稀疏实数阵）、\fld{build\_id}（原子递增缓存键）；
  \item \textbf{投影回溯}：\fld{bus\_merge\_map}、\fld{projection\_certificate}、
        \fld{projection\_report}（\texttt{std::optional}）。
\end{itemize}
配套映射类型 \srcpath{BusIndexMap}/\srcpath{BranchIndexMap}/\srcpath{SystemIndexMap}
见 \srcpath{include/hacdcpf/assembly/index\_map.hpp:BusIndexMap}。

\subsection{装配流程}
入口 \srcpath{make\_solver\_data(const HybridPowerSystem\&, LossModelType)}
（\srcpath{src/power\_flow/solver\_data.cpp:make\_solver\_data}），步骤：
\begin{enumerate}
  \item rich$\to$canonical 投影 \srcpath{projection::RichToCanonicalOperator::apply}
        （:268--269，含开关/零阻抗合并）；
  \item \srcpath{materialize\_dc\_storage} 与逐字段 move；\texttt{StaticGeneratorDC}
        转换为 \texttt{StaticGenerator}（:make\_solver\_data）；
  \item \fld{build\_id} 取自原子计数器（:266,319），作为全部缓存的失效标记；
  \item \textbf{ExternalGrid 处理}（:make\_solver\_data）：把母线 \fld{vm\_pu}/\fld{va\_deg}
        写为外部电网整定，并将该母线 PQ 提升为 SLACK；
  \item \textbf{成网换流器提升} \srcpath{apply\_acpv\_voltage\_control}（:apply\_acpv\_voltage\_control）：
        按 \srcpath{resolve\_device\_control\_role}，AC\_GRID\_FORMING 换流器的交流母线
        提升为 SLACK 并钉住 \fld{v\_ac\_set\_pu}/\fld{v\_ac\_angle\_set\_deg}；
        AC\_PV / DC\_V\_DROOP\_AC\_V 模式提升为 PV；已有 SLACK 不动；
  \item \srcpath{rebuild\_matrices}（:rebuild\_matrices）$=$ 发电聚合 $+$ 负荷聚合 $+$
        $Y_{bus}$ 构建 $+$ $G_{dc}$ 构建。
\end{enumerate}

\subsection{注入聚合（指定功率）}
\subsubsection*{发电聚合}
\srcpath{aggregate\_generation}（src/power\_flow/solver\_data.cpp:aggregate\_generation）把下列注入聚进
\fld{pg}/\fld{qg}（除以 \fld{base\_mva} 转 pu）：发电机、
\texttt{StaticGenerator}（$p\cdot$\fld{scaling}）、可再生能源、
PV（经功率曲线 \srcpath{compute\_pv\_power\_mw}，见第~\ref{sec:distribution}~章）、
储能（放电为正）、VPP（PCC 交换功率）、并网微网（\fld{p\_exchange} 出口为正）、
非在途移动储能。

\subsubsection*{负荷聚合与 ZIP 归一化}
\srcpath{aggregate\_load\_demand}（src/power\_flow/solver\_data.cpp:aggregate\_load\_demand）：
Load 表为空（且无充电站）时 \fld{has\_component\_loads}=false，回退母线级
\fld{pd\_mw}/\fld{qd\_mvar}；否则以母线级需求为底（恒定功率，ZIP 种子 $(1,0,0)$），
叠加 $p_{\mathrm{mw}}\cdot$\fld{scaling} 的分量负荷，并按 $|P|/|Q|$ 加权平均
归一化出每母线 ZIP 系数（:160--168, 193--211）。充电站恒功率叠加
（kW$\to$MW，:173--188）。潮流方程中的指定注入为
（\srcpath{src/power\_flow/jacobian\_builder.cpp:build\_power\_spec}）：
\begin{align*}
P_i^{\mathrm{spec}} &= P_{g,i}-p_{d,i}\big(w_{p0}+w_{p1}V_i+w_{p2}V_i^2\big)\\
Q_i^{\mathrm{spec}} &= Q_{g,i}-q_{d,i}\big(w_{q0}+w_{q1}V_i+w_{q2}V_i^2\big)
\end{align*}
权重顺序 $w_0$=恒功率、$w_1$=恒电流、$w_2$=恒阻抗。VSC 交流注入经
\srcpath{converter\_ac\_injection} 加入 spec（src/power\_flow/jacobian\_builder.cpp:build\_power\_spec）。

\subsection{交流导纳阵 $Y_{bus}$}
\srcpath{build\_admittance\_matrix(data)}（\srcpath{src/power\_flow/admittance\_builder.cpp:build\_admittance\_matrix}），
π 型支路 + 复变比（MATPOWER \texttt{makeYbus} 约定）。对在役支路 $i\to j$
（\fld{from\_bus}/\fld{to\_bus} 为 1-based，代码中 $-1$）：
\begin{align*}
y_s &= \frac{1}{r+\mathrm{j}x},\qquad
y_{tt} = y_s+\mathrm{j}\frac{b}{2},\qquad
\hat{t} = |\mathrm{tap}|\,\mathrm{e}^{\mathrm{j}\,\mathrm{shift}\cdot\pi/180}\\
Y_{ff} &= \frac{y_{tt}}{|\hat{t}|^2},\qquad
Y_{ft} = -\frac{y_s}{\hat{t}^{*}},\qquad
Y_{tf} = -\frac{y_s}{\hat{t}}
\end{align*}
对角元累加 $Y_{ii}\leftarrow Y_{ii}+Y_{ff}$、$Y_{jj}\leftarrow Y_{jj}+y_{tt}$，
非对角元 $Y_{ij}=Y_{ft}$、$Y_{ji}=Y_{tf}$（:build\_admittance\_matrix）。装配层为 fail-close 口径：
$|z|$ 非有限或 $\le\texttt{kMinSeriesImpedancePu}=10^{-12}$ 即抛
\texttt{std::runtime\_error}（:50--55，阈值 :15）；\fld{tap} 非有限或
$\le 10^{-12}$ 同样抛出（:build\_admittance\_matrix）——零阻抗支路须在投影层先行合并，
装配层不做静默回退。母线对地并联（$V=1.0$~pu 功率口径，
MATPOWER 惯例）直接入对角（:build\_admittance\_matrix）：
\[
Y_{ii}\leftarrow Y_{ii}+\frac{g_s+\mathrm{j}b_s}{S_{\mathrm{base}}}
\]
独立 \texttt{Shunt} 表同样入对角；可投切并联按
\fld{bs\_per\_step}$\times$\fld{current\_step} 取值（:build\_admittance\_matrix）。
装配用三元组列表 + \texttt{setFromTriplets}。

\subsection{直流电导阵 $G_{dc}$}
\srcpath{build\_dc\_conductance(data)}（src/power\_flow/admittance\_builder.cpp:build\_dc\_conductance）：
正对角 Laplacian。对在运 DC 支路取（$r_{pu}$ 非有限或 $\le 10^{-12}$ 即抛
\texttt{std::runtime\_error}，fail-close，:128--132）
\[
g=\frac{1}{r_{pu}},\qquad G_{ii}\leftarrow G_{ii}+g,\qquad G_{ij}=G_{ji}=-g
\]
（:build\_dc\_conductance）。DC 节点功率方程为 $\mathbf p_{dc}=\mathbf v_{dc}\odot(G_{dc}\mathbf v_{dc})$
（src/power\_flow/jacobian\_builder.cpp:evaluate\_residual\_impl）。

\subsection{FDPF 电纳阵 $B'/B''$}
\srcpath{build\_susceptance\_matrix}（src/power\_flow/admittance\_builder.cpp:build\_susceptance\_matrix）取
$-\mathrm{Im}(Y^{\mathrm{mod}}_{ij})$，构型由 \texttt{BpBuildFlags} 控制
（\fld{zero\_resistance}/\fld{zero\_charging}/\fld{unit\_tap}/\fld{zero\_phase\_shift}/%
\fld{add\_bus\_shunts}/\fld{min\_x\_pu}，头文件
\srcpath{include/hacdcpf/assembly/ybus\_builder.hpp:BpBuildFlags}）：
\begin{itemize}
  \item $B'$（\srcpath{make\_bp\_flags}）：\textbf{FDXB} 方案——$R=0$、忽略充电、
        变比幅值取 1、相移取 0、不含母线并联；$|x|<$\fld{min\_x\_pu} 的支路跳过；
  \item $B''$（\srcpath{make\_bpp\_flags}）：\textbf{FDBX} 风格——保留 $R$、充电与变比
        （即 $-\mathrm{Im}(Y)$ 全量）、相移取 0，并加母线并联 $-B_s/S_{\mathrm{base}}$。
\end{itemize}

\subsection{DC 注入装配与残差独立复算}
\srcpath{assemble\_dc\_injections(data, vm, va, vdc, pdc\_spec)}
（\srcpath{src/power\_flow/pf\_injection\_assembly.cpp:assemble\_dc\_injections}），净注入约定发电为正：
DC 负荷为负（Load 表空时回退母线 \fld{pd\_mw}，:25--37）、DC 储能放电为正、
静态发电机 $p\cdot$\fld{scaling}、PV 阵列 \fld{p\_set}、VSC 经
\srcpath{converter\_dc\_injection}、DC/DC（入端 $-P_{in}$、出端 $+P_{out}$，
含 $I^2R$ 损耗，:79--88）、能量路由器 DC 端口。DC 参考选择
\srcpath{first\_dc\_slack\_or\_default}（\srcpath{include/hacdcpf/power\_flow/pf\_utils.hpp:first\_dc\_slack\_or\_default}）。

\srcpath{evaluate\_power\_flow\_residual}（\srcpath{src/power\_flow/residual\_evaluator.cpp:evaluate\_power\_flow\_residual}）
提供独立于 Newton 热路径的残差复算，返回 \{ac, dc, full\} 三块 $\ell_\infty$ 范数，
供诊断与测试交叉校验。交流注入采用\textbf{稀疏复数形式}
\srcpath{compute\_ac\_power\_injections}（:compute\_ac\_power\_injections）：$\mathbf V=\mathrm{polar}(V_m,\theta)$，
$\mathbf I=Y_{bus}\mathbf V$ 为一次稀疏矩阵-向量乘，$S_i=V_i\,\overline{I_i}$
（:23--32，调用点 :72），复杂度 $O(\mathrm{nnz}(Y_{bus}))$，不再把 $Y_{bus}$
稠密化后做 $O(n^2)$ 双重循环（数学上与逐元素 $g\cos\theta+b\sin\theta$ 求和等价）。
其母线分类用
\srcpath{classify\_ac\_buses}（\srcpath{include/hacdcpf/power\_flow/pf\_utils.hpp:classify\_ac\_buses}，额外 SLACK 归 PV），
与 Newton 上下文（额外 SLACK 整体剔除）略有差异。

\subsection{支路功率回算}
求解收敛后 \srcpath{compute\_branch\_flows}（\srcpath{src/power\_flow/branch\_flow.cpp:compute\_branch\_flows}）
按 π 等值 + 复变比逐支路回算两端功率（:compute\_branch\_flows）：
\begin{align*}
S_f &= V_f\,\overline{\big(y_{ff}V_f+y_{ft}V_t\big)}\cdot S_{\mathrm{base}}\\
S_t &= V_t\,\overline{\big(y_{tf}V_f+y_{tt}V_t\big)}\cdot S_{\mathrm{base}}
\end{align*}
输出 \texttt{BranchFlow\{pf\_mw, qf\_mvar, pt\_mw, qt\_mvar\}}（MW/Mvar），
按 \fld{data.ac\_branches} 顺序；停运/零阻抗支路保持 0。注意：本文件名为
"branch\_flow" 但\textbf{不是} Baran--Wu 支路潮流法；配网前推回代见
第~\ref{sec:distribution}~章。

\subsection{验证与校核}
装配层的验证以端到端方式覆盖：85 个 MATPOWER 算例的电压对照直接检验 $Y_{bus}$
装配（含分接头/移相/充电的正确性）；\srcpath{tools/transformer\_tap\_crosscheck.py}
对 12 组双节点变压器配置做 pandapower/MATPOWER/本平台三向对照；ZIP 聚合与
换流器提升逻辑由 \srcpath{tests/test\_all\_components\_pf.cpp}（29 类组件逐类禁用、
解变化量阈值断言）与 \srcpath{tests/test\_dc\_island\_reference.cpp} 覆盖；
投影幂等性、合并映射复合与结果恢复由投影蜕变测试与
\srcpath{tools/gui\_api\_e2e.py}（Bus 1 P/Q 闭合）保证。
